\section{Introduction}
\label{sec:intro}

The gravitational-wave signal of a compact-binary merger fixes its luminosity
distance without any rung of the astronomical distance ladder, which makes
every detection an absolute standard siren \citep{Schutz1986, HolzHughes2005}.
Converting that distance into a measurement of the expansion rate requires a
redshift, and the one merger with an identified electromagnetic counterpart,
GW170817, delivered \Hz\ to about ten per cent \citep{Abbott2017SirenH0}.
Counterparts have not repeated. Dark standard sirens supply the missing
redshift statistically instead: rather than identifying a single host, the
analysis uses the redshift distribution of all catalogued galaxies consistent
with the event's sky localisation and marginalises over which one was the host
\citep{DelPozzo2012}. Applied to the events detected so far, that method
constrains \Hz\ to a few tens of per cent without a single counterpart
\citep{Fishbach2019, Palmese2023, GWTC3Cosmo2023, GWTC4Cosmo2025}. The
motivation for pushing it harder is external: the cosmic microwave background
gives $67.4 \pm 0.5$~\kmsmpc\ \citep{Planck2018} while the Cepheid distance
ladder gives $73.0 \pm 1.0$~\kmsmpc\ \citep{Riess2022}, and a route to \Hz\
that shares systematics with neither will arbitrate between them only once it
reaches the per cent level.

The precision of a dark-siren inference is set by the host catalog as much as
by the gravitational-wave data. Two catalog properties do the work. The first
is how sharply the catalog pins the redshift of a host given the event's sky
patch, which improves with catalog density and with the amount of structure in
the redshift distribution on the scale of the localisation
\citep{Fishbach2019, Gray2020}. The second is how well the galaxies the survey
missed are accounted for, since real catalogs are flux limited and beyond some
redshift most hosts are absent
\citep{ChenFishbachHolz2018, Finke2021, Gray2022}. Both favour the densest
available catalog, and dark-siren analyses have accordingly been built on
galaxy catalogs alone \citep{Palmese2023, GWTC3Cosmo2023}.

There is a second host population the data have something to say about. Several
formation channels place a subset of binary black hole mergers in the gas-rich,
dense environments around AGN, where gas torques and migration traps harden
binaries and can assemble massive systems through repeated mergers
\citep{Bartos2017, StoneMetzgerHaiman2017, McKernan2018,
TagawaHaimanKocsis2020}. Individual candidate counterparts have made the
association plausible without establishing it
\citep{Graham2020, Gayathri2021, BomPalmese2023}. The fraction \fagn\ of
mergers supplied by that channel is the branching ratio between AGN-assisted
assembly and everything else \citep{FordMcKernan2022}. What is known about it
is known indirectly: \fagn\ is read off features of the mass and spin
distributions that the channel is expected to imprint, at a precision of a
factor of a few \citep{McKernan2018, FordMcKernan2022}.

AGN are not distributed like galaxies, and that is what makes \fagn\
measurable from the same events that measure the expansion rate. A catalog of
luminous AGN is roughly two orders of magnitude sparser in comoving number
density than a galaxy catalog reaching the same depth \citep{Ananna2022}, and
AGN sit in more massive haloes, so they are more strongly clustered, with a
larger linear bias \citep{Krumpe2010}. Both differences enter the
line-of-sight redshift prior each catalog supplies, so the two priors differ in
normalisation and in shape, the sparse one resolving individual structures
along a line of sight where the dense one is nearly smooth. The weight of the
mixture between them is therefore identifiable. Host positions have been used
to test the association from the outside, by asking whether the detected
events' localisation volumes contain more AGN of a given class than chance
allows \citep{Bartos2017NatComm, Veronesi2022, Veronesi2023}. No dark-siren
analysis, however, has carried a second host population inside its redshift
prior, where the same association becomes a parameter measured jointly with
the expansion rate rather than a hypothesis tested against it.

We measure \Hz\ and \fagn\ jointly on a simulated universe in which both are
known and in which the two catalogs carry the densities and biases the
literature reports. One clustered density field is sampled by two tracers, a
dense galaxy population and an AGN population two orders of magnitude sparser;
\EvNobs\ mergers are drawn from the two host populations in a known proportion
and observed once each, with every measurement error a function of the observed
data rather than of the source; and the same catalog is thinned by a magnitude
limit into flux-limited versions, so that the measurement can be repeated on
the catalogs a survey would deliver.

Section~\ref{sec:methods} gives the likelihood: the mixture over catalogs, the
redshift prior each catalog supplies and how its sky dependence is normalised,
the treatment of flux-limited catalogs, and the selection function.
Section~\ref{sec:data} describes the simulated universe: one clustered density
field sampled by two tracers with literature-anchored densities and biases, and
the events observed against it. Section~\ref{sec:results} reports the
measurements: \Hz\ from one catalog at a time, \Hz\ and \fagn\ jointly, each catalog on its own hosts at equal event
counts, and both parameters on flux-limited catalogs.
Section~\ref{sec:validation} tests them, by repeating the entire construction
on independent realisations of the simulated universe and by a null test that
removes the host association while changing nothing else.
Section~\ref{sec:discussion} states what the measurement establishes and what
it depends on. Uncertainties are 68\% equal-tailed credible intervals unless
stated otherwise, a value quoted as $x \pm y$ carries its standard error, and
the input cosmology is flat $\Lambda$CDM with
$\Hz = \HzeroTruth$~\kmsmpc\ and $\Omega_{m,0} = \OmZeroTruth$.
